\chapter{Narcissistic Personality Disorder}

\medskip
\noindent\textit{\textbf{Under review.} This chapter is an AI-assisted educational draft; it is not a substitute for professional medical advice.}

\section*{Definition and Overview}
Narcissistic personality disorder is a persistent pattern of grandiosity, need for admiration, and reduced empathy that causes significant distress or impairment. Occasional self-focus or confidence does not establish the disorder, and diagnosis requires a comprehensive assessment.

\section*{Clinical Features}
Possible features include an exaggerated sense of importance, preoccupation with success or status, entitlement, exploitation of relationships, sensitivity to criticism, shame, unstable self-esteem, and difficulty recognising other people's needs. Presentation varies.

\section*{Diagnosis and Differential Diagnosis}
Diagnosis is clinical and considers long-term patterns across settings. Assessment distinguishes the disorder from normal personality variation, bipolar disorder, depression, trauma-related symptoms, substance effects, and other personality disorders.

\section*{Management}
Psychotherapy is the main treatment and may focus on emotional regulation, relationships, self-esteem, empathy, and coping with criticism. Medicines may treat coexisting depression, anxiety, or other symptoms but do not specifically cure the personality pattern.

\section*{When to Seek Help}
Seek professional help when relationship, work, mood, substance-use, or safety problems are persistent. Call 000 or attend an emergency department for immediate danger, serious self-harm, or inability to remain safe.

\section*{Expanded clinical context}
Narcissistic personality disorder NPD is a mental-health topic. Interpretation should account for age, baseline health, medicines, comorbidities, goals, and access to care. This chapter supports discussion with a qualified clinician and does not establish a diagnosis.

\section*{Assessment and decision points}
Assessment should consider onset, duration, triggers, sleep, substance use, physical health, developmental context, social stressors, and immediate safety. Ask about self-harm or harm from others when appropriate. Document the main concern, time course, functional impact, relevant risks, and red flags. Use targeted investigations when their result is likely to change management.

\section*{Management and follow-up}
Management may include education, psychological therapy, practical support, treatment of coexisting conditions, medication when indicated, and an agreed follow-up or safety plan. Explain expected improvement, when reassessment is needed, and how follow-up can be accessed. Avoid starting, stopping, or sharing prescription medicines without professional advice.

\section*{Safety-netting}
Seek urgent help for immediate danger, suicidal intent, psychosis, severe confusion, inability to meet basic needs, or rapidly worsening symptoms.

\section*{Practical prevention}
Use plain-language information, supportive relationships, regular routines, and professional review when symptoms persist or impair daily life. Keep written instructions and seek review if the topic recurs, changes, or does not improve as expected.

\section*{Limitations}
This educational expansion should be checked against current local guidance and authoritative topic-specific sources.
\section*{A simple starting plan}
Use a consistent routine, identify situations that trigger conflict or distress, and discuss practical goals with a clinician. Progress is usually assessed by functioning, relationships, and emotional regulation rather than by a single symptom.

\clearpage
